\documentclass[../../main.tex]{subfiles}

\begin{document}

\section{参考文献}\label{section:ch10-参考文献}

\S10.1：此处给出的粘性解定义归功于 Crandall、Evans 和 Lions（Trans. AMS 282 (1984), 487--502），他们重新表述了 Crandall 和 Lions 在基础论文（Trans. AMS 277 (1983), 1--42）中提出的更早定义。Tran~\cite{Tn} 是该主题的良好导引。

\S10.2：唯一性证明遵循上述第一篇论文，改进之处我得自 M.~Crandall。

\S10.3：P.-L.~Lions~\cite{Li} 首先观察到粘性解定义与控制论最优性条件之间的联系。Fleming--Soner~\cite{F-S} 和 Bardi--Capuzzo~Dolcetta~\cite{B-CD} 的书提供了更多关于粘性解及其与确定性和随机最优控制联系的信息。

\S10.4：D.~Ostrov 贡献了\cref{exercise:ch10-从明暗恢复形状}。关于完全非线性椭圆型方程的更多内容参见 Cabr\'e--Caffarelli~\cite{C-C} 和 Gilbarg--Trudinger~\cite{G-T}。

\end{document}